% GENERATED FILE. Edit canonical lessons, then run npm run generate:books.
% canonical-source-hash: fnv1a64:1d2363b8a3bd3393
% canonical-lessons: ZH-C541-yaobushi, ZH-C541-duokui, ZH-C541-gaishan, ZH-C541-zhuyao, ZH-C541-yaome, ZH-R541-review-537-541, ZH-R541-review-final-541

\chapter{If It Weren't for You}
\label{ch:zh-if-it-weren-t-for-you}
\hlchaptermodality{541}

\begin{chapteropening}
\textbf{By the end of this chapter:} \emph{I can say if it weren't for, thanks to, things will improve, what the main reason is, and either this or that.}

Complete the last lesson of chapter 541: I can say if it weren't for, thanks to, things will improve, what the main reason is, and either this or that.
\end{chapteropening}

\section[\texorpdfstring{\zh{要不是} (yàobúshì)}{要不是 (yàobúshì)}]{\zh{要不是} (yàobúshì) — if it were not for}

\label{lesson:ZH-C541-yaobushi}



\begin{quote}
Before the new one: say the Mandarin for the whole chain of causes, then the Mandarin for what is going on here.
\end{quote}

\subsection*{You'll want to know: \zh{要不是}}
\textbf{\zh{要不是}} — \emph{yàobúshì} — \textquotedblleft{}if it were not for\textquotedblright{}.

\begin{quote}
\textbf{\zh{要不是你我们就找不到这儿了}} \emph{yàobúshì nǐ wǒmen jiù zhǎo bu dào zhèr le} — If it weren't for you, we'd never have found this place.
\end{quote}

Tone: \zh{不} is said bú before the falling tone of shì, and the spelling shows it.

\begin{cousinweb}
\textbf{\zh{要不是}} (\emph{yàobúshì}) is built from \zh{要} 'if', \zh{不} 'not' and \zh{是} 'to be', all inherited from Middle Chinese, \emph{'jiewH}, \emph{pjut} and \emph{dzyeX} (Guangyun, in Baxter's transcription); the dictionary glosses the first 'to bind by agreement', and its use for 'if' is not given there.
\end{cousinweb}

\begin{grammarlens}[title={What you've built}]
One more word for saying what would happen.
\end{grammarlens}

\subsection*{Your turn}
\begin{itemize}
\raggedright
  \item \emph{Say it:} \emph{yàobúshì}
  \item \emph{Say it:} \emph{yàobúshì}, once more
  \item \emph{Say it:} \emph{yàobúshì nǐ wǒmen jiù zhǎo bu dào zhèr le} — If it weren't for you, we'd never have found this place.
  \item \emph{Recall:} say the Mandarin for the whole chain of causes, then the Mandarin for what is going on here, then say \emph{yàobúshì} again
\end{itemize}

\begin{tcolorbox}[breakable,colback=teal!4,colframe=teal!35!black,title={Before you move on}]
What does \zh{要不是} mean? (\textquotedblleft{}If it were not for\textquotedblright{}.) Say it once more.
\end{tcolorbox}
\section[duōkuī]{duōkuī — it is thanks to; luckily for}

\label{lesson:ZH-C541-duokui}



\begin{quote}
Before the new one: say the Mandarin for what is going on here, then the Mandarin for if it were not for.
\end{quote}

\subsection*{You'll want to know: duōkuī}
\emph{duōkuī} — \textquotedblleft{}it is thanks to; luckily for\textquotedblright{}.

\begin{quote}
\emph{duōkuī nǐ bāng wǒ bùrán wǒ jiù zuò bu wán} — Thanks to your help I finished; otherwise I'd never have managed.
\end{quote}

Its characters come later; for now it is a word for the ear.

\begin{cousinweb}
\emph{duōkuī} is built from \zh{多} 'much' and \emph{kuī}, both inherited from Middle Chinese, \emph{ta} and \emph{khjwe} (Guangyun, in Baxter's transcription), which glosses the first 'many' and the second 'to be lacking'. How the pair came to mean 'thanks to' is not settled.
\end{cousinweb}

\begin{grammarlens}[title={What you've built}]
One more word for saying what would happen.
\end{grammarlens}

\subsection*{Your turn}
\begin{itemize}
\raggedright
  \item \emph{Say it:} \emph{duōkuī}
  \item \emph{Say it:} \emph{duōkuī}, once more
  \item \emph{Say it:} \emph{duōkuī nǐ bāng wǒ bùrán wǒ jiù zuò bu wán} — Thanks to your help I finished; otherwise I'd never have managed.
  \item \emph{Recall:} say the Mandarin for what is going on here, then the Mandarin for if it were not for, then say \emph{duōkuī} again
\end{itemize}

\begin{tcolorbox}[breakable,colback=teal!4,colframe=teal!35!black,title={Before you move on}]
What does duōkuī mean? (\textquotedblleft{}It is thanks to; luckily for\textquotedblright{}.) Say it once more.
\end{tcolorbox}
\section[gǎishàn]{gǎishàn — to make conditions better}

\label{lesson:ZH-C541-gaishan}



\begin{quote}
Before the new one: say the Mandarin for if it were not for, then the Mandarin for it is thanks to.
\end{quote}

\subsection*{You'll want to know: gǎishàn}
\emph{gǎishàn} — \textquotedblleft{}to make conditions better\textquotedblright{}.

\begin{quote}
\emph{zhǐyào dàjiā yìqǐ nǔlì huánjìng jiù huì gǎishàn} — As long as we all work together, the environment will get better.
\end{quote}

Its characters come later; for now it is a word for the ear.

\begin{cousinweb}
\emph{gǎishàn} is built from \zh{改} 'to change' and \emph{shàn} 'good', both inherited from Middle Chinese, \emph{kojX} and \emph{dzyenX} (Guangyun, in Baxter's transcription), which glosses the first 'to alter' and the second 'good'.
\end{cousinweb}

\begin{grammarlens}[title={What you've built}]
One more word for saying what would happen.
\end{grammarlens}

\subsection*{Your turn}
\begin{itemize}
\raggedright
  \item \emph{Say it:} \emph{gǎishàn}
  \item \emph{Say it:} \emph{gǎishàn}, once more
  \item \emph{Say it:} \emph{zhǐyào dàjiā yìqǐ nǔlì huánjìng jiù huì gǎishàn} — As long as we all work together, the environment will get better.
  \item \emph{Recall:} say the Mandarin for if it were not for, then the Mandarin for it is thanks to, then say \emph{gǎishàn} again
\end{itemize}

\begin{tcolorbox}[breakable,colback=teal!4,colframe=teal!35!black,title={Before you move on}]
What does gǎishàn mean? (\textquotedblleft{}To make conditions better\textquotedblright{}.) Say it once more.
\end{tcolorbox}
\section[zhǔyào]{zhǔyào — main; chief}

\label{lesson:ZH-C541-zhuyao}



\begin{quote}
Before the new one: say the Mandarin for it is thanks to, then the Mandarin for to make conditions better.
\end{quote}

\subsection*{You'll want to know: zhǔyào}
\emph{zhǔyào} — \textquotedblleft{}main; chief\textquotedblright{}.

\begin{quote}
\emph{wǒ bú qù de zhǔyào yuányīn shì tiānqì bù hǎo} — The main reason I'm not going is the bad weather.
\end{quote}

Its characters come later; for now it is a word for the ear.

\begin{cousinweb}
\emph{zhǔyào} is built from \emph{zhǔ} 'to be in charge' and \zh{要} 'important', both inherited from Middle Chinese, \emph{tsyuX} and \emph{'jiewH} (Guangyun, in Baxter's transcription), which glosses the first 'to manage, to lead' and the second 'to bind by agreement'.
\end{cousinweb}

\begin{grammarlens}[title={What you've built}]
One more word for saying what would happen.
\end{grammarlens}

\subsection*{Your turn}
\begin{itemize}
\raggedright
  \item \emph{Say it:} \emph{zhǔyào}
  \item \emph{Say it:} \emph{zhǔyào}, once more
  \item \emph{Say it:} \emph{wǒ bú qù de zhǔyào yuányīn shì tiānqì bù hǎo} — The main reason I'm not going is the bad weather.
  \item \emph{Recall:} say the Mandarin for it is thanks to, then the Mandarin for to make conditions better, then say \emph{zhǔyào} again
\end{itemize}

\begin{tcolorbox}[breakable,colback=teal!4,colframe=teal!35!black,title={Before you move on}]
What does zhǔyào mean? (\textquotedblleft{}Main; chief\textquotedblright{}.) Say it once more.
\end{tcolorbox}
\section[\texorpdfstring{\zh{要么} (yàome)}{要么 (yàome)}]{\zh{要么} (yàome) — either one or the other; either … or}

\label{lesson:ZH-C541-yaome}



\begin{quote}
Before the new one: say the Mandarin for to make conditions better, then the Mandarin for the main one.
\end{quote}

\subsection*{You'll want to know: \zh{要么}}
\textbf{\zh{要么}} — \emph{yàome} — \textquotedblleft{}either one or the other; either … or\textquotedblright{}.

\begin{quote}
\textbf{\zh{要么你来要么我去}} \emph{yàome nǐ lái yàome wǒ qù} — Either you come here or I go there.
\end{quote}

\begin{cousinweb}
\textbf{\zh{要么}} (\emph{yàome}) is built from \zh{要} and the light ending \zh{么}, as in \zh{怎么}. The first is inherited from Middle Chinese \emph{'jiewH} (Guangyun, in Baxter's transcription); the second is said light, and its history in this use is not settled.
\end{cousinweb}

\begin{grammarlens}[title={What you've built}]
One more word for saying what would happen.
\end{grammarlens}

\subsection*{Your turn}
\begin{itemize}
\raggedright
  \item \emph{Say it:} \emph{yàome}
  \item \emph{Say it:} \emph{yàome}, once more
  \item \emph{Say it:} \emph{yàome nǐ lái yàome wǒ qù} — Either you come here or I go there.
  \item \emph{Recall:} say the Mandarin for to make conditions better, then the Mandarin for the main one, then say \emph{yàome} again
\end{itemize}

\begin{tcolorbox}[breakable,colback=teal!4,colframe=teal!35!black,title={Before you move on}]
What does \zh{要么} mean? (\textquotedblleft{}Either one or the other; either … or\textquotedblright{}.) Say it once more.
\end{tcolorbox}
\section[(dialogue)]{Looking back: from As It Turns Out to If It Weren't for You}

\label{lesson:ZH-R541-review-537-541}



\begin{quote}
Say the Mandarin for what set it off and as it turns out.
\end{quote}

\subsection*{Your turn}
\begin{itemize}
\raggedright
  \item \emph{Say it:} what set it off, as it turns out, to result in something bad, to cause damage, how well it works
  \item \emph{Say it:} so as not to, and thereby, in itself, to arise from something, not in the least
  \item \emph{Say it:} to stir up trouble, to put the blame on, a strange business, to come off worse, to get into trouble
  \item \emph{Say it:} eccentric, to drag someone into trouble, bound to happen, the whole chain of causes, what is going on here
  \item \emph{Say it:} if it were not for, it is thanks to, to make conditions better, the main one, either one or the other
  \item \emph{Read it:} \zh{原} \zh{怪} — each one aloud, then a word that holds it
\end{itemize}

\begin{tcolorbox}[breakable,colback=teal!4,colframe=teal!35!black,title={Before you move on}]
What set it off? (\textbf{\zh{起因}}, \emph{qǐyīn}.) As it turns out? (\textbf{\zh{原来}}, \emph{yuánlái}.) To result in something bad? (\emph{dǎozhì}.) To cause damage? (\emph{zàochéng}.) How well it works? (\emph{xiàoguǒ}.) So as not to? (\emph{miǎnde}.) And thereby? (\emph{cóng'ér}.) In itself? (\textbf{\zh{本身}}, \emph{běnshēn}.) To arise from something? (\emph{chǎnshēng}.) Not in the least? (\emph{gēnběn}.) To stir up trouble? (\emph{rě}.) To put the blame on? (\textbf{\zh{怪}}, \emph{guài}.) A strange business? (\textbf{\zh{怪事}}, \emph{guàishì}.) To come off worse? (\emph{chīkuī}.) To get into trouble? (\emph{chuǎnghuò}.) Eccentric? (\textbf{\zh{古怪}}, \emph{gǔguài}.) To drag someone into trouble? (\emph{liánlei}.) Bound to happen? (\emph{nánmiǎn}.) The whole chain of causes? (\textbf{\zh{前因后果}}, \emph{qiányīn hòuguǒ}.) What is going on here? (\textbf{\zh{怎么回事}}, \emph{zěnme huí shì}.) If it were not for? (\textbf{\zh{要不是}}, \emph{yàobúshì}.) It is thanks to? (\emph{duōkuī}.) To make conditions better? (\emph{gǎishàn}.) The main one? (\emph{zhǔyào}.) Either one or the other? (\textbf{\zh{要么}}, \emph{yàome}.)
\end{tcolorbox}
\section[(dialogue)]{If It Weren't for You — once more}

\label{lesson:ZH-R541-review-final-541}



\begin{quote}
Say the Mandarin for if it were not for and it is thanks to.
\end{quote}

\subsection*{Your turn}
\begin{itemize}
\raggedright
  \item \emph{Say it:} if it were not for, it is thanks to, to make conditions better, the main one, either one or the other
\end{itemize}

\begin{tcolorbox}[breakable,colback=teal!4,colframe=teal!35!black,title={Before you move on}]
If it were not for? (\textbf{\zh{要不是}}, \emph{yàobúshì}.) It is thanks to? (\emph{duōkuī}.) To make conditions better? (\emph{gǎishàn}.) The main one? (\emph{zhǔyào}.) Either one or the other? (\textbf{\zh{要么}}, \emph{yàome}.)
\end{tcolorbox}
